% 摘要
%
% 要求：150-250 词。结构建议：
%   1. 本综述考察什么（1-2 句）
%   2. 我们提出了什么分类体系（1-2 句）
%   3. 我们得出了什么判断（2-3 句）
%   4. 留下什么开放问题（1 句）
%
% ⚠️ 摘要必须与 06-conclusion.tex 的核心结论一致。
% ⚠️ pdflatex 无法输出中文：正文与 \todo{} 内容一律用英文，
%    中文说明只能写在以 % 开头的注释行里。

\begin{abstract}
Recursive self-improvement (RSI) requires that an improved system participate
in producing the next improvement. This review examines the outermost of three
commonly distinguished loci of self-improvement: the \emph{harness} --- the
program structure that organizes a foundation model's reasoning, tool use, and
action --- as distinct from model parameters or training data. We survey
\todo{N} primary studies spanning \todo{2023 to 2026}, comprising \todo{N}
peer-reviewed conference papers, \todo{N} industrial reports, and \todo{N}
preprints from 2025 and 2026.

We organize this literature along five axes: the \emph{target} of improvement
(prompt, control flow, source code, tools, memory, evaluator, or context
strategy), the \emph{signal} driving improvement, the \emph{search operator},
the \emph{strength of guarantees}, and --- most consequentially --- the
\emph{depth of recursion}. We find that \todo{one-sentence core claim, in
English}.

Several claims in the current literature remain unverified, particularly
regarding whether reported gains survive scrutiny for selection bias, compute
cost, and the independence of the verifier from the system being optimized.
\todo{Close with the two or three open problems named in Section 5.}
\end{abstract}
